\scanpage{107}
\uline{Corollaire} Soit $\varphi$ \unsure{repr.|réal.} quelconque de \deleted{\unsure{$\underline G$}}
$\Pi$ polygone combinatoire de degré $N$
($3\leq N\leq+\infty$). \inserted{Soit $\underline u$ un gén de $\mathrm{Aut}^+(\Pi)$, $u$ son image en $\mathrm{Aff}(E)$.} Alors les entiers suivants
sont tous égaux :

\[
\text{1)}\quad\mathrm{Card}\,(S=\mathrm{Im}(\underline S\unsure{-|\to}E))
\]

\[
\text{2)}\quad\mathrm{Card}\,(A=\mathrm{Im}(\underline A\to\mathrm{Dr}(E))
\]

\[
\text{3)}\quad\mathrm{Card}\,(G^+=\mathrm{Im}\,\mathrm{Aut}^{+}(\Pi)\to\mathrm{Aff}(E))
\]

\[
\text{4) Ordre de }u\in\mathrm{Aff}(E)
\]

\deleted{\unsure{Par}} Cet entier \inserted{$n$} est \deleted{$\geq3$ et est} un diviseur de $N$
(condition vide si $N=+\infty$) est $\unsure{\infty|\text{ou}}\geq3$, et c'est
aussi le plus \deleted{grand} \inserted{petit des} diviseurs \inserted{$n$} de $N$ tels
que $\varphi$ se factorise par le \inserted{polygone} quotient
\inserted{$\Pi_{(n)}$} de $\Pi$ d'ordre $n$, et enfin l'unique
diviseur $n$ de $N$ tel que $\varphi$ se factorise
par le polygone quotient $\Pi_{(n)}$ de $\Pi$
d'ordre $n$ \uline{de façon} que la réalisation
obtenue de $\Pi_{(n)}$ soit \uline{fidèle}.

\begin{center}\rule{1.6cm}{0.35pt}\end{center}

Mais il manque un \unsure{sorite|sorte} sur les
morphismes \inserted{(étales)} de polygônes \inserted{comb.} (ou plus
gén\textsuperscript{t} de \unsure{cartes|contours} \inserted{comb.}\ldots), nous permettant
% archive sheet number: 106
